% Pista ana-24i-q5 · Anàlisi
% Procedència: PAU Catalunya 2024 · Convocatòria de juny (incidències) · Sèrie 5 · Qüestió 5
\documentclass[11pt,a4paper]{article}
\usepackage{pista}
\pistainfo{Catalunya 2024}{Convocatòria de juny (incidències)}{Sèrie 5}

\begin{document}

\section*{\textcolor{blauPista}{Qüestió 5}}

\noindent En Carles vol construir un decorat per a l'obra de teatre de final de curs en forma d'un rectangle i dos semicercles. Les variables $x$ i $y$ són les semi-mides indicades a la figura.

\subsection*{\textit{a)} \normalfont Determineu el perímetre i l'àrea del decorat que s'ha de construir en funció de $x$ i de $y$. \hfill\textnormal{[1 punt]}}

\pista{Identifica clarament, a partir de la figura, què és cadascuna de les parts que componen el decorat: quines són les dimensions del rectangle i quins són els radis dels dos semicercles? El perímetre del decorat és la suma de les longituds dels arcs dels dos semicercles (cadascun és mitja circumferència) més els trams de la vora del rectangle que no queden ``tapats'' per cap semicercle. L'àrea és la suma de l'àrea del rectangle més les dues àrees dels semicercles.}

\subsection*{\textit{b)} \normalfont Per a revestir el perímetre del decorat, en Carles té material per a cobrir fins a $10 \text{ m}$. Si el vol gastar tot, quines seran les mides del decorat d'àrea màxima que podrà construir? Quin és el valor d'aquesta àrea? \hfill\textnormal{[1,5 punts]}}

\pista{Imposa la condició $P = 10$ (perímetre igual a 10 m) i aïlla una de les dues variables en funció de l'altra. Substitueix aquesta expressió a la fórmula de l'àrea: quedarà una funció d'una sola variable. A partir d'aquí és un problema d'optimització: deriva, resol l'equació derivada igual a zero, i comprova que el valor trobat és efectivament un màxim. Ves amb compte de respondre TOTES les parts de la pregunta: les \emph{mides} (totes dues) i l'\emph{àrea} màxima.}

\end{document}
